% ============================================================================
% COURS 09 — MACHINE À COURANT CONTINU
% Reprise fidèle de 13-MCC.docx
% ============================================================================

\section{Généralités sur les machines électriques}\label{sec:mcc-generalites}

\subsection{Notions d'électromagnétisme}

Un courant électrique crée un champ magnétique. Réciproquement, une variation de flux
magnétique peut induire une force électromotrice. Les aimants possèdent deux pôles : nord et
sud. Deux pôles identiques se repoussent, tandis que deux pôles opposés s'attirent.

Le champ d'induction magnétique est noté $\vec B$ et s'exprime en teslas. Dans les machines
électriques, les inductions usuelles sont de l'ordre de $1$ à $1{,}5\,\mathrm T$.

\subsection{Principe simplifié des machines électriques}

Le fonctionnement d'une machine tournante peut être interprété comme l'interaction entre un
champ statorique $\vec B_s$ et un champ rotorique $\vec B_r$. Ces champs sont produits par des
enroulements parcourus par des courants ou par des aimants permanents.

On distingue notamment :
\begin{itemize}
  \item les machines à champ statorique fixe, comme la machine à courant continu ;
  \item les machines à champ statorique tournant, comme les machines synchrones et asynchrones.
\end{itemize}

\section{Constitution de la machine à courant continu}\label{sec:mcc-constitution}

La MCC comporte :
\begin{itemize}
  \item un inducteur ou stator, source du champ magnétique ;
  \item un induit ou rotor, siège des forces électromotrices ;
  \item un circuit magnétique comprenant stator, rotor et entrefer ;
  \item un ensemble balais--collecteur assurant l'alimentation du rotor et la commutation des courants.
\end{itemize}

\subsection{Inducteur}

L'inducteur peut être constitué d'aimants permanents ou d'un enroulement d'excitation. Dans le
second cas, le flux $\Phi$ est commandé par le courant d'excitation $I_e$, tant que le circuit
magnétique n'est pas saturé.

\subsection{Induit et collecteur}

L'induit est constitué d'un circuit magnétique feuilleté et d'un bobinage. Le collecteur est
solidaire du rotor ; ses lames de cuivre sont reliées aux enroulements. Les balais sont fixes et
assurent la liaison électrique avec l'induit.

Le système balais--collecteur redresse les forces électromotrices élémentaires et inverse le
courant dans les conducteurs au moment opportun afin de maintenir le couple dans un même sens.

\subsection{Modes d'excitation}

Les principaux modes sont : excitation indépendante, excitation shunt, excitation série et
excitation composée. Dans les petites puissances, les aimants permanents sont très fréquents.

\section{Relations fondamentales et bilan de puissance}\label{sec:mcc-relations}

\subsection{Force électromotrice}

La force électromotrice de l'induit est proportionnelle au flux et à la vitesse :
\[
E=K\Phi\Omega.
\]
À flux constant :
\[
E=K_e\Omega,
\]
où $K_e$ est la constante de force électromotrice en $\mathrm{V\,s\,rad^{-1}}$.

\subsection{Modèle de l'induit en régime établi}

\TLSchemaInduitMCCPermanent

En convention récepteur :
\[
U=E+RI.
\]
En multipliant par $I$ :
\[
UI=EI+RI^2.
\]
On identifie la puissance absorbée $P_a=UI$, les pertes Joule $P_J=RI^2$ et la puissance
électromagnétique $P_{em}=EI$.

\subsection{Couple électromagnétique et couple utile}

La puissance électromagnétique vérifie :
\[
P_{em}=EI=C_{em}\Omega.
\]
À flux constant :
\[
C_{em}=K_c I.
\]
En unités SI, les constantes $K_e$ et $K_c$ ont la même valeur numérique.

Le couple utile disponible sur l'arbre est :
\[
C_u=C_{em}-C_p,
\]
où $C_p$ regroupe les pertes mécaniques et ferromagnétiques.

Lors d'un essai à vide, $C_u=0$, donc $C_{em}=C_p$. On peut ainsi déterminer le couple de
pertes à partir du courant à vide :
\[
C_p=K_c I_0.
\]

\subsection{Commande à flux constant}

À flux constant, deux propriétés rendent la MCC simple à commander :
\begin{itemize}
  \item le courant d'induit impose directement le couple électromagnétique ;
  \item la tension d'induit impose principalement la vitesse.
\end{itemize}

\TLQuadrantsMCC

\subsection{Bilan de puissance et rendement}

\TLBilanPuissancesMCC

Le rendement moteur vaut :
\[
\eta=\frac{P_u}{P_a}=\frac{C_u\Omega}{UI}.
\]

\section{Fonctionnement en régime variable}\label{sec:mcc-variable}

Lorsque le courant varie, l'inductance de l'induit doit être prise en compte.

\TLSchemaInduitMCCVariable

Les quatre équations fondamentales à flux constant sont :
\[
\boxed{U=RI+L\frac{\mathrm dI}{\mathrm dt}+K_e\Omega}
\]
\[
\boxed{C_{em}=K_c I}
\]
\[
\boxed{J\frac{\mathrm d\Omega}{\mathrm dt}=C_{em}-C_r-f\Omega}
\]
\[
\boxed{E=K_e\Omega}
\]

Dans le domaine de Laplace, le modèle peut être représenté sous forme de schéma-blocs :

\TLSchemaBlocMCC

Le couple résistant apparaît comme une perturbation mécanique.

\section{Fonctionnement en génératrice}\label{sec:mcc-generatrice}

Lorsque la charge entraîne la machine, la puissance mécanique est convertie en puissance
électrique. La machine fonctionne alors en génératrice. Le signe du produit $C\Omega$ devient
négatif.

En convention générateur :
\[
U=E-RI
\]
en régime établi. Pour récupérer l'énergie, la chaîne électrique, le variateur et la source
d'alimentation doivent être réversibles.

\section{Démarrage de la MCC}\label{sec:mcc-demarrage}

Au démarrage, $\Omega=0$, donc $E=0$. Le courant initial vaut alors :
\[
I_d=\frac{U}{R}.
\]
Comme la résistance d'induit est faible, ce courant peut être très élevé. Le couple de démarrage
est :
\[
C_d=K_c I_d.
\]
Pour les puissances importantes, on démarre sous tension réduite ou avec une limitation de
courant. On limite généralement l'appel de courant entre $1{,}25I_n$ et $2I_n$.

\section{Variation de vitesse}\label{sec:mcc-variation}

En régime établi et à flux constant :
\[
\Omega=\frac{U-RI}{K_e}.
\]
Trois actions sont possibles.

\subsection{Variation de la résistance d'induit}

Ajouter une résistance série réduit la vitesse mais augmente les pertes Joule. Cette méthode est
peu utilisée.

\subsection{Variation de la tension d'induit}

\TLVariationVitesseMCC

C'est la solution la plus courante. La variation de tension est réalisée par un convertisseur
statique : hacheur pour une source continue, pont commandé pour une source alternative.

\subsection{Variation du flux inducteur}

La diminution du flux permet de dépasser la vitesse nominale. Cette zone de défluxage se fait à
puissance approximativement constante, donc au prix d'une diminution du couple disponible.

\section{Identification expérimentale du modèle}\label{sec:mcc-identification}

\TLEssaisIdentificationMCC

\subsection{Identification de $K_e$}

À vide, le courant est faible et la chute $RI$ peut être négligée. En régime permanent :
\[
U\simeq E=K_e\Omega.
\]
La pente de la droite $U=f(\Omega)$ fournit $K_e$.

\subsection{Identification de $R$ et $L$}

Rotor bloqué, $E=0$. Sous un échelon de tension réduit :
\[
U=RI+L\frac{\mathrm dI}{\mathrm dt}.
\]
La réponse est du premier ordre :
\[
I(t)=\frac{U}{R}\left(1-e^{-t/\tau}\right),\qquad \tau=\frac{L}{R}.
\]
La valeur finale fournit $R$ et la constante de temps fournit $L$.

\section{Applications et exercices intégrés}\label{sec:mcc-applications}

Le document source propose plusieurs applications : borne solaire, vélo à assistance électrique,
système de treillis, magasin vertical, scooter électrique, béquille électrique, kart et maison
hantée. Elles mobilisent les mêmes démarches :
\begin{enumerate}
  \item établir le schéma équivalent ;
  \item écrire les équations électriques et mécaniques ;
  \item identifier les paramètres par essais ;
  \item calculer les puissances et le rendement ;
  \item déterminer les quadrants et la réversibilité nécessaires.
\end{enumerate}

\section{À retenir}\label{sec:mcc-retenir}

\begin{itemize}
  \item $E=K_e\Omega$ à flux constant.
  \item $C_{em}=K_cI$ et, en SI, $K_e=K_c$ numériquement.
  \item En régime établi moteur : $U=E+RI$.
  \item En régime variable : $U=E+RI+L\,\mathrm dI/\mathrm dt$.
  \item Le courant commande le couple ; la tension commande principalement la vitesse.
  \item La MCC est naturellement réversible, mais la récupération dépend du variateur et de la source.
  \item Le démarrage doit être limité car $E=0$ lorsque $\Omega=0$.
\end{itemize}
